\section{Описание практической части}
\label{sec:implementation}

Раздел описывает практическую реализацию программного комплекса:
средства разработки, архитектуру, ключевые модули и их интерфейсы.

\subsection{Выбор средств разработки}

В качестве среды разработки выбраны следующие программные
компоненты, обеспечивающие необходимую воспроизводимость и
модульность системы.

\begin{itemize}
    \item \textbf{Robot Operating System 2 (ROS~2), дистрибутив Humble
          Hawksbill.} Фреймворк обеспечивает стандартизованную модель
          распределённого программного обеспечения робототехнической
          системы: узлы, топики, действия, сервисы и дерево
          преобразований систем координат.
    \item \textbf{Gazebo Classic.} Физически-реалистичный симулятор
          рабочей среды, поддерживающий плагины омни-движения
          (\texttt{gazebo\_ros\_planar\_move}) и модели сенсоров.
    \item \textbf{Navigation~2 (Nav2).} Навигационный стек ROS~2,
          используемый для серверов глобального планировщика,
          локального контроллера, поведений и BT~Navigator.
    \item \textbf{slam\_toolbox.} Пакет online-SLAM на базе 2D-лидара.
    \item \textbf{OpenCV} и \texttt{cv\_bridge}~--- для обработки
          изображений и цветовой сегментации.
    \item \textbf{Docker, Docker~Compose, tmux, tmuxp}~--- для
          контейнеризации среды исполнения и оркестрации стартовой
          последовательности.
    \item \textbf{Xvfb, x11vnc, noVNC}~--- для удалённого доступа к
          графическим интерфейсам симулятора и RViz независимо от
          графической подсистемы хоста.
\end{itemize}

\subsection{Состав узлов и их назначение}

Ключевые программные узлы и их функции приведены в
таблице~\ref{tab:nodes}.

\begin{table}[H]
\caption{Ключевые узлы программного комплекса}\label{tab:nodes}
\centering
\small
\begin{tabularx}{\textwidth}{|>{\raggedright\arraybackslash}p{0.30\textwidth}|>{\raggedright\arraybackslash}X|}
\hline
\textbf{Узел} & \textbf{Функция} \\
\hline
\texttt{obstacles\_controller} & Создание цели и цилиндрических препятствий через сервис \texttt{/spawn\_entity}, периодическая публикация случайных скоростей в топики \texttt{/\textless{}name\textgreater/cmd\_vel} для имитации динамической среды. \\
\hline
\texttt{target\_detector} & Детекция синего объекта в кадре камеры в пространстве HSV, оценка дальности до цели по данным 2D-лидара, публикация флага видимости \texttt{/target\_visible} и позы \texttt{/target\_pose} в системе координат карты. \\
\hline
\texttt{target\_nav\_bridge} & Формирование и отправка целей \texttt{NavigateToPose} в Nav2 по данным детектора, запуск поведения \texttt{Spin} при потере цели для её поиска. \\
\hline
\texttt{scan\_bridge} & Ретрансляция топика \texttt{/omni\_robot/scan} в \texttt{/scan} для совместимости с настройками SLAM и costmap. \\
\hline
\texttt{cmd\_vel\_bridge} & Ретрансляция управляющих скоростей Nav2 из \texttt{/cmd\_vel} в топик плагина движения робота \texttt{/omni\_robot/cmd\_vel}. \\
\hline
\texttt{odom\_tf\_bridge} & Публикация преобразования \texttt{odom} $\to$ \texttt{base\_link} по данным одометрии плагина движения. \\
\hline
\texttt{controller\_server} (Nav2) & Загрузка плагина \texttt{MPPIController} с настроенным набором критиков, вычисление команд скорости на каждом такте. \\
\hline
\texttt{slam\_toolbox} & Построение карты помещения в режиме online, публикация карты и преобразования \texttt{map} $\to$ \texttt{odom}. \\
\hline
\end{tabularx}
\end{table}

\subsection{Модель робота}

Модель робота описана в формате URDF/xacro и включает:
\begin{itemize}
    \item базовое тело \texttt{base\_link} в виде
          параллелепипеда с задаваемыми габаритами для аппроксимации корпуса робота (по умолчанию
          0{,}3$\times$0{,}3$\times$0{,}1~м);
    \item 2D-лидар на верхней грани корпуса с углом обзора $360\degree{}$,
          360 лучами и диапазоном 0{,}1--10~м (плагин
          \texttt{gazebo\_ros\_ray\_sensor});
    \item фронтальную RGB-камеру с разрешением
          640$\times$480 и горизонтальным полем зрения 1{,}047~рад
          (плагин \texttt{gazebo\_ros\_camera});
    \item плагин движения \texttt{libgazebo\_ros\_planar\_move.so},
          принимающий на вход сообщение \texttt{geometry\_msgs/Twist} с
          компонентами $v_x$, $v_y$ и $\omega$, что соответствует
          модели~\eqref{eq:kin_matrix}.
\end{itemize}

Параметры модели вынесены в описание робота и меняются без изменения
логики контроллера, что соответствует принципу параметризации,
сформулированному в требованиях.

\subsection{Среда симуляции и подвижные сущности}

В качестве базовой среды симуляции используется мир квартиры
\texttt{apartment.world}, построенный на основе открытой модели
AWS~RoboMaker (small\_house). Мир содержит стены, мебель и свободное
пространство, достаточное для маневрирования робота.

Подвижные сущности создаются пакетом \texttt{apartment\_sim} в виде
цилиндрических моделей с плагином движения
\texttt{libgazebo\_ros\_planar\_move.so}. Узел
\texttt{obstacles\_controller} публикует в их топики \texttt{cmd\_vel}
случайные скорости, формируя сценарий с движением цели и
препятствий в ограниченной области. Параметры (число препятствий,
скорости, места появления, геометрия цилиндров) задаются через
ROS-параметры.

\reportfig[0.95]{house.png}{Сцена симуляции квартиры}{fig:house_scene_1}
\reportfig[0.95]{house_up.png}{Сцена симуляции квартиры (вид сверху)}{fig:house_scene_2}

\subsection{Контур восприятия цели}

Узел \texttt{target\_detector} реализован на Python с использованием
OpenCV и \texttt{cv\_bridge}. На каждом кадре камеры выполняются
следующие шаги:
\begin{enumerate}
    \item преобразование кадра из пространства BGR в HSV;
    \item бинаризация маски по диапазону оттенка синего
          (\texttt{hsv\_blue\_lower}, \texttt{hsv\_blue\_upper}) с
          дополнительными критериями доминирования синего канала над
          красным и зелёным;
    \item морфологическое открытие с заданным размером ядра;
    \item выделение связных компонент и выбор наибольшего контура;
    \item оценка горизонтального углового направления на цель по
          нормированной координате центра цели и горизонтальному полю
          зрения камеры;
    \item оценка дальности до цели по данным лидара в окне индексов
          вокруг полученного направления (медиана допустимых значений);
    \item пересчёт полученной точки в координатную систему
          \texttt{map} через TF.
\end{enumerate}

Опубликованные топики:
\begin{itemize}
    \item \texttt{/target\_visible} (std\_msgs/Bool)~--- флаг видимости
          цели;
    \item \texttt{/target\_pose} (geometry\_msgs/PoseStamped)~--- поза
          цели в карте;
    \item \texttt{/target\_marker} (visualization\_msgs/Marker)~--- маркер
          для визуализации в RViz;
    \item \texttt{/omni\_robot/camera/image\_raw/target\_status}~---
          аннотированный кадр с индикатором видимости.
\end{itemize}

Такой набор выходов позволяет модулю MPPI и верхнему уровню логики
использовать только высокоуровневые признаки цели, без привязки к
деталям обработки изображения.

\subsection{Формирование навигационных целей}

Узел \texttt{target\_nav\_bridge} связывает контур восприятия с
действием Nav2 \texttt{NavigateToPose}. Его логика:
\begin{enumerate}
    \item при наличии актуальной позы цели и отличии её от последней
          отправленной более чем на заданный порог~--- обновить цель
          Nav2;
    \item при переходе \texttt{/target\_visible} из \texttt{True} в
          \texttt{False}~--- сначала позволить контуру довести движение
          до последней известной позы, затем запустить режим
          \texttt{Spin} (вращение на месте) для поиска цели;
    \item при возврате видимости~--- отменить режим \texttt{Spin}
          и вернуться в режим сопровождения.
\end{enumerate}

Такой конечный автомат уменьшает число резких переключений цели и
стабилизирует поведение при временной потере наблюдаемости.

\subsection{Локальное управление с MPPI}

Локальным контроллером в стеке Nav2 служит плагин
\texttt{nav2\_mppi\_controller::MPPIController}, работающий в режиме
\texttt{motion\_model = Omni}, что соответствует модели~\eqref{eq:kin_matrix}.
Основные параметры конфигурации (файл
\texttt{config/omni\_nav2\_params.yaml}):
\begin{itemize}
    \item горизонт прогнозирования: $N = 56$ шагов;
    \item шаг дискретизации: $\Delta t = 0{,}05$~с;
    \item число сэмплов: $M = 2000$;
    \item стандартные отклонения возмущений по
          $v_x, v_y, \omega$: $0{,}24$, $0{,}24$, $0{,}8$ соответственно;
    \item верхние/нижние границы скоростей:
          $v_x \in [-0{,}65; 0{,}95]$~м/с, $v_y \in [-0{,}95; 0{,}95]$~м/с,
          $\omega \in [-3{,}8; 3{,}8]$~рад/с;
    \item температура $\lambda = 0{,}3$;
    \item частота работы контроллера 50~Гц.
\end{itemize}

Применён следующий набор критиков, соответствующий слагаемым
функции~\eqref{eq:stage_cost}:
\begin{itemize}
    \item \texttt{ConstraintCritic}~--- штраф за нарушение ограничений
          скорости~\eqref{eq:vel_limits};
    \item \texttt{CostCritic}, \texttt{ObstaclesCritic}~--- изотропный штраф
          за приближение к препятствиям на основе costmap, реализация
          $\Psi_{\text{obs},k}$~\eqref{eq:psi_obs};
    \item \texttt{PotentialFieldCritic}~--- направленный штраф
          $\Psi_{\text{pf},k}$~\eqref{eq:psi_pf}, реализованный в
          рамках ВКР;
    \item \texttt{GoalCritic}~--- штраф за отклонение конечной позы от
          цели (терминальное слагаемое);
    \item \texttt{TargetCameraCritic}~--- пользовательский критик
          центрирования цели в поле зрения камеры, реализующий
          слагаемое $w_{\phi}\,(e^{\phi}_k)^{2}$ из~\eqref{eq:stage_cost};
    \item \texttt{PathAlignCritic}, \texttt{PathFollowCritic}~--- штрафы,
          удерживающие траекторию рядом с путём глобального планировщика;
    \item \texttt{TwirlingCritic}~--- штраф за избыточное вращение.
\end{itemize}

Критики \texttt{GoalAngleCritic} и \texttt{PathAngleCritic} в данной
конфигурации отключены: для омни-платформы желаемая ориентация в
финальной позе и угол сближения с путём не существенны, а
центрирование цели обеспечивается \texttt{TargetCameraCritic}.

\subsubsection{Пользовательский критик \texttt{TargetCameraCritic}}

В методе \texttt{score} критика:
\begin{enumerate}
    \item получается последняя опубликованная поза цели
          \texttt{/target\_pose} и проверяется её свежесть по таймауту;
    \item поза цели пересчитывается в кадр траекторий контроллера
          (глобальная карта) через TF;
    \item для каждой (batch, timestep) пары вычисляется желаемый
          азимут на цель
          ${\phi^{\star}}^{(m)}_{k+j} = \operatorname{atan2}\bigl(y^{t} - y_{k+j}^{(m)},\,
          x^{t} - x_{k+j}^{(m)}\bigr)$;
    \item формируется средняя по траектории абсолютная ошибка
          $|e^{\phi}_k|$ между желаемым азимутом и углом рыскания
          корпуса;
    \item получившееся значение, умноженное на настроечный вес, в
          степени \texttt{cost\_power} добавляется в суммарную стоимость
          траекторий MPPI.
\end{enumerate}

Критик автоматически отключается при устаревании топики
\texttt{/target\_pose}, что согласовано с логикой узла
\texttt{target\_nav\_bridge} и поддерживает корректное поведение при
потере цели.

\subsubsection{Пользовательский критик \texttt{PotentialFieldCritic}}

Критик реализует направленный штраф
$\Psi_{\text{pf},k}$~\eqref{eq:psi_pf}, восполняющий ограничение
изотропных штрафов \texttt{CostCritic} и \texttt{ObstaclesCritic} для
голономных движений. На каждом такте для текущей позы робота и для
точек сэмплированных траекторий выполняется:
\begin{enumerate}
    \item чтение значений costmap в точке и вокруг неё (радиус
          \texttt{activation\_sample\_radius}), если максимум ниже
          порога активации \texttt{activation\_cost\_threshold}~--- критик
          возвращает нулевую стоимость;
    \item вычисление расстояния до препятствия $\tilde{d}(c)$ по
          обратной характеристике инфляционного слоя costmap
          (параметры \texttt{inflation\_radius},
          \texttt{cost\_scaling\_factor}~--- те же, что у
          \texttt{InflationLayer});
    \item оценка двухточечного градиента $\nabla c$ по двум
          разносам \texttt{gradient\_sample\_distance} (с весами 0{,}7
          и~0{,}3) для устойчивости к численному шуму costmap и
          построение направления отталкивания
          $\hat{\mathbf{f}}$~\eqref{eq:pf_dir};
    \item суммирование вдоль траектории до прохождения дистанции
          \texttt{influence\_distance}: проксимальная составляющая
          $w_{\rho}\,\rho_k$ и направленная составляющая
          $w_{\sigma}\,(1+\rho_k)\,\sigma_k$ из~\eqref{eq:pf_prox},
          \eqref{eq:pf_dir_score};
    \item ограничение полученной стоимости снизу нулём и сверху
          параметром \texttt{max\_force\_cost} с последующим
          возведением в степень \texttt{cost\_power} и суммированием
          с весом \texttt{cost\_weight}, помноженным на масштаб
          \texttt{dominance\_scale}.
\end{enumerate}

Внутри радиуса \texttt{near\_goal\_distance} от текущей цели Nav2
критик выключается, чтобы не препятствовать сближению с финальной
позой. Для отладки в топик \texttt{/potential\_field\_critic/markers}
публикуется визуализация направления $\hat{\mathbf{f}}$ из текущей
позы робота, что позволяет проверять корректность оценки градиента в
RViz.

\subsection{SLAM}

Для построения карты и локализации используется \texttt{slam\_toolbox}
в асинхронном режиме. Параметры (файл \texttt{config/slam.yaml}):
решатель Ceres с линейным методом
\texttt{SPARSE\_NORMAL\_CHOLESKY} и стратегией Levenberg--Marquardt,
разрешение карты 0{,}05~м, максимальная дальность лидара 15~м,
интервал обновления карты 1~с.

Навигационный стек использует \texttt{map} как глобальный фрейм (для
планировщика и глобальной costmap) и \texttt{odom} как локальный фрейм
(для локальной costmap и локального контроллера), что соответствует
рекомендациям Nav2 по разделению ответственности.

\subsection{Пример рабочей сцены}

Пример рабочего состояния системы в процессе сопровождения цели
приведён на рисунке~\ref{fig:rviz}. На снимке экрана RViz видно:
\begin{itemize}
      \item красная стрелка - направление действия PotentialFieldCritic,
      \item зеленая стрелка - ориентация робота,
      \item оранжевая стрелка - направление на цель,
      \item построенную карту помещения,
      \item глобальную costmap,
      \item показания 2D-лидара,
      \item глобальный путь и маркер цели,
      \item кадр фронтальной камеры с наложенным индикатором видимости цели.
\end{itemize}

\reportfig[0.95]{rviz screenshot.png}{Рабочая сцена в RViz: карта помещения, costmap, лидар, путь и маркер цели}{fig:rviz}